% =====================================================================
%  TYPE: Drug prescription / medication schedule
%  Standalone version of the last page of the postop instructions:
%  follow-up visit, ER visit, non-surgical patient. Does NOT replace
%  the e-prescription (ΗΔΙΚΑ) for reimbursed drugs.
%
%  Copy, rename to "YYYYMMDD - surname - topic.tex", fill in and
%  compile with XeLaTeX. Toggle alternatives by moving the "%".
%  Text in [brackets] is a placeholder: replace or delete it.
% =====================================================================
\documentclass[10pt,greek]{article}
\usepackage[useregional]{datetime2}

%%%%%%%%%%%%%%%%%%%%% Language
\providecommand\DocLang{gr}

%%%%%%%%%%%%%%%%%%%%% Clinic (letterhead logo, footer address)
\providecommand\Clinic{Peristeri}
% \providecommand\Clinic{Mandra}
% \providecommand\Clinic{Bioclinic}
% \providecommand\Clinic{Psychikou}
% \providecommand\Clinic{Metropolitan}
% \providecommand\Clinic{Dunant}
% \providecommand\Clinic{Euroclinic}

%%%%%%%%%%%%%%%%%%%%% Document code for the QR: paste the code issued
%%%%%%%%%%%%%%%%%%%%% by doc.erns.eu (empty = no QR code)
\providecommand\DocCode{}

%%%%%%%%%%%%%%%%%%%%% Patient
\providecommand\PatientPrefix{Ο}
% \providecommand\PatientPrefix{Η κα}
\providecommand\PatientName{Παπαδόπουλος Γιώργος}
\providecommand\PatientAge{45}
% \providecommand\PatientAge{}

%%%%%%%%%%%%%%%%%%%%% Date
\providecommand\DocDate{\DTMdate{2026-01-15}}
% \providecommand\DocDate{\DTMdate{2026-05-04}}

%%%%%%%%%%%%%%%%%%%%% Title
\providecommand\DocTitle{Οδηγίες Φαρμακευτικής Αγωγής}
% \providecommand\DocTitle{Συνταγή}

%%%%%%%%%%%%%%%%%%%%% Closing (empty = no signature block)
\providecommand\DocClosing{Ο ιατρός}
% \providecommand\DocClosing{}

%%%%%%%%%%%%%%%%%%%%% Page numbers: all = also on page 1 when the prescription has several pages
\providecommand\DocPageNumbers{all}
% \providecommand\DocPageNumbers{yes}   % from page 2 on only
% \providecommand\DocPageNumbers{no}

%%%%%%%%%%%%%%%%%%%%% Stamp and signature image
\providecommand\DocStamp{no}
% \providecommand\DocStamp{yes}

\input{evrtemplate}

\begin{document}

\begin{PatientInfo}
%   \Info{Διάγνωση}{[διάγνωση]}
%   \Info{Αλλεργίες}{[καμία γνωστή]}
%   \Info{Ημ. Χειρουργείου}{DD-MM-YYYY}
\end{PatientInfo}

\begin{MedTable}
\DrugEntry{tb Paracetamol}{1000mg}{1 x 3}{10}{}
% ---- Drug library: uncomment the lines you need ---------------------
% -- Analgesics
% \DrugEntry{tb Paracetamol}{1000mg}{1 x 3}{10}{}
% \DrugEntry{tb Apotel}{1000mg}{2 ή 3 την ημέρα}{}{παυσίπονο, μέγιστη δόση 3 δισκία το 24ωρο}
% \DrugEntry{tb Depon}{500mg}{1 ή 2, έως 3 φορές την ημέρα}{}{εφόσον υπάρχει πόνος -- μέγιστο 3g το 24ωρο}
% \DrugEntry{tb Tramal}{50mg}{1 x 2 ή 3}{6}{}
% \DrugEntry{tb Zaldiar}{}{2 ή 3 την ημέρα}{}{εφόσον υπάρχει έντονος πόνος}
% -- NSAIDs
% \DrugEntry{tb Movatec}{7.5mg}{1 x 2}{10}{μετά το φαγητό}
% \DrugEntry{tb Movatec}{15mg}{1 x 1}{1}{μετά το φαγητό}
% \DrugEntry{tb Naprosyn}{500mg}{1 x 2}{6}{μετά το φαγητό}
% -- Steroids (tapers)
\DrugEntry{tb Medrol}{4mg}{1 x 3\newline 1 x 2\newline 1 x 1\newline STOP}{}{για 5 ημέρες\newline για 5 ημέρες\newline για 5 ημέρες}
% \DrugEntry{tb Medrol}{4mg}{2 x 3\newline 1 x 3\newline 1 x 2\newline 1 x 1\newline STOP}{}{για 3 ημέρες\newline για 3 ημέρες\newline για 3 ημέρες\newline για 3 ημέρες}
% \DrugEntry{syr Dexaton}{2mg/5ml}{10cc x 3}{10}{μετά το φαγητό}
% \DrugEntry{tb Dexamethasone}{4mg}{1 x 3}{}{μετά το φαγητό}
% -- Antibiotics
\DrugEntry{tb Zinadol}{500mg}{1 x 2}{5}{αντιβίωση}
% \DrugEntry{tb Klaricid}{500mg}{1 x 2}{7}{αντιβίωση}
% \DrugEntry{caps Amoxicillin}{500mg}{1 x 3}{7}{μετά το φαγητό}
% -- Gastroprotection
\DrugEntry{caps Omeprazole}{40mg}{1 x 2}{10}{προ φαγητού}
% \DrugEntry{caps Omeprazole}{20mg}{1 x 2}{10}{15 λεπτά πριν το φαγητό}
% \DrugEntry{caps Mostrelan}{40mg}{1 x 1}{10}{προ φαγητού}
% -- Anticoagulants
% \DrugEntry{inj Clexane}{4000IU}{1 x 1}{10}{αντιπηκτικές ενέσεις υποδορίως, την ίδια ώρα κάθε ημέρα}
% \DrugEntry{inj Innohep}{0,45ml}{1 x 1}{10}{αντιπηκτικές ενέσεις υποδορίως, την ίδια ώρα κάθε ημέρα}
% \DrugEntry{inj Fraxiparine}{2850IU}{1 x 1}{10}{αντιπηκτικές ενέσεις υποδορίως, την ίδια ώρα κάθε ημέρα}
% -- Neuropathic pain, antiepileptics
% \DrugEntry{tb Lyrica}{50mg}{1 x 1 (βράδυ)}{4}{μετά την 4η ημέρα, 1 x 2}
% \DrugEntry{caps Lyrica}{150mg}{1 x 2}{}{πρωί -- βράδυ, με τροφή. Μην διακόπτεται απότομα}
% \DrugEntry{tb Neurontin}{400mg}{2 x 2}{}{}
\DrugEntry{tb Tegretol}{200mg}{1 x 3}{}{με τροφή, σταθερές ώρες. \textbf{Μην διακόπτεται απότομα}}
% \DrugEntry{caps Cymbalta}{60mg}{1 x 2}{}{}
% \DrugEntry{tb Keppra}{1000mg}{1 x 3}{}{}
% \DrugEntry{tb Depakine}{400mg}{1 x 3}{}{θα γίνει προοδευτική μείωση σε συνεννόηση με τον θεράποντα ιατρό}
% -- Other
% \DrugEntry{sir Duphalac}{5ml}{1 ή 2 την ημέρα}{}{καθαρτικό, ανάλογα με τη λειτουργία του εντέρου -- στόχος: μαλακές κενώσεις χωρίς προσπάθεια}
% \DrugEntry{tb Xanax}{0,25mg}{1 x 2-3}{}{μέχρι 3 την ημέρα επί εκνευρισμού}
% \DrugEntry{tb Escitalopram}{10mg}{1 x 1}{}{}
% \DrugEntry{spray Vioplex-T ή Pulvo47}{}{1 x 2}{}{στο τραύμα, και να καλύπτεται μετά}
% ---------------------------------------------------------------------
\end{MedTable}

% Notes under the table (uncomment the ones you need)
\textbf{Σημείωση:} Μην λαμβάνετε αντιφλεγμονώδη φάρμακα (π.χ.\ Voltaren, Brufen, Ponstan, Naprosyn) χωρίς ιατρική οδηγία.
% \textbf{Σημείωση:} Μην διακόψετε και μην μειώσετε μόνος/μόνη σας τα [Tegretol / Lyrica / Keppra]. Οι δόσεις θα μειωθούν σταδιακά κατόπιν ιατρικής οδηγίας.
% \textbf{Σημείωση:} Ενημερώστε κάθε ιατρό και φαρμακοποιό ότι λαμβάνετε [φάρμακο], πριν από τη χορήγηση οποιουδήποτε νέου φαρμάκου.
% \textbf{Σημείωση:} Η υπόλοιπη φαρμακευτική σας αγωγή (π.χ.\ αντιυπερτασική) συνεχίζεται \textbf{αμετάβλητη}.

\end{document}
